%%
%% fall-lecture05.tex
%% 
%% Made by Alex Nelson <pqnelson@gmail.com>
%% Login   <alex@lisp>
%% 
%% Started on  2026-10-08T10:29:13-0700
%% Last update 2026-10-08T10:29:13-0700
%% 

\lecture{}

\begin{definition}
Let $(X,\topology)$ be a topological space. Let $U\subset X$ and
$G\subset X$ be arbitrary subsets such that
$U\subset G$. Then we write $U\Subset G$ if $\closure{U}\subset G$.
\end{definition}

\begin{lemma}[Urysohn]
Let $(X,\topology)$ be a normal space. Let $F_{0}$ and $F_{1}$ be two
disjoint subsets of $X$. Then there exists a continuous function
$f\colon X\to[0,1]$ such that $f|_{F_{0}}=0$ and $f|_{F_{1}}=1$.
\end{lemma}

\begin{proof}[Proof (for metric spaces)]
Set
\begin{equation}
f(x)=\frac{\rho(x,F_{0})}{\rho(x,F_{0})+\rho(x,F_{1})},
\end{equation}
which is continuous, the denominator is never zero, and takes values
in $[0,1]$. This satisfies the lemma.
\end{proof}

\begin{proof}[Proof (arbitrary spaces)]
Consider the following situation: we have a closed set $F$, and an
open set $G$ and $F\subset G$. Then $G^{\complement}=X\setminus G$
is closed. Then we use the $T_{4}$ axiom to find a neighborhood $U$ of
$FR$ and a neighborhood $V$ of $G^{\complement}$. This means we have
by separation $F\Subset U\Subset G$. We will form a family $U_{q}$ of
subsets where $q\in\QQ\cap(0,1)$. We define $U_{1/2}:=U$.

Now, we iterate this construction using $F=F_{0}$ and $G=U_{1/2}$ to
get $U_{1/3}$; and similarly for $F=U_{1/2}$ and $G=G$ we obtain
$U_{2/3}$. After one step we get
\begin{equation}
F_{0}\Subset U_{1/3}\Subset U_{1/2}\Subset U_{2/3}\Subset G.
\end{equation}
The book-keeping with the indices may be understood as follows. The
first step, we have $0/1$ and $1/1$, then obtain
\begin{equation}
\frac{0+1}{1+1}=\frac{1}{2},
\end{equation}
which is why we named $U_{1/2}:=U$. The second step, we have the
indices $(0/1, 1/2, 1/1)$ and we obtain
\begin{equation}
\frac{0+1}{1+2}=\frac{1}{3},\quad\mbox{and}\quad\frac{1+1}{2+1}=\frac{2}{3}.
\end{equation}
The third step $(0/1, 1/3)$ gives $1/4$, $(1/3,1/2)$ gives us $2/5$,
$(1/2,2/3)$ gives us $3/5$, $(2/3,1/1)$ gives us $3/4$. And so on
\textit{ad infinitum}. This sequence of indices are the \define{Farey sequence}
of rational numbers written in reduced form.

We claim: given disjoint closed subsets $F_{0}$, $F_{1}$, there exists
a family $U_{q}$ of neighborhoods where $q\in\QQ\cap(0,1)$ and for any
$r<s$ we have $F_{0}\Subset U_{r}\Subset U_{s}\Subset F_{1}$. (We
should think of this like a topographic map $U_{r}$ has height $r\in\QQ\cap(0,1)$.)

Define
\begin{equation}
f(x):=\inf\{r\in\QQ\cap(0,1)\mid x\in U_{r}\}.
\end{equation}
We claim
\begin{enumerate}
\item $f=0$ on $F_{0}$ [Pf: let $x\in F_{0}$ be arbitrary, then $x\in U_{r}$
  for all $r\in\QQ\cap(0,1)$. Therefore $f(x)=\inf(\QQ\cap(0,1))=0$.]
\item $f=1$ on $F_{1}$ [Pf: similar to (1)]
\item $f$ is continuous.
\end{enumerate}
Now, for continuity, look at the preimages of open sets. The preimage
of $f^{-1}((-\infty,\alpha))=\{x\in X\mid f(x)<\alpha\}$. But
$f(x)=\alpha$ iff $\inf\{r\mid x\in U_{r}\}=\alpha$ iff $x\in U_{r}$
for some $r<\alpha$ iff $x\in\bigcup U_{r}$ which is open. Similarly,
$f^{-1}((\alpha,+\infty))=\bigcup_{r>\alpha}(\closure{U_{r}})^{\complement}$
which is open.

Lemma 1: Half-lines $(-\infty, a)$ and $(b,+\infty)$ where $a,b\in\QQ$
generate the standard topology on $\RR$.

Lemma 2: If $f\colon(X,\topology_{X})\to(Y,\topology_{Y})$ is a map
between topological spaces and $\mathcal{E}\subset\powerset{Y}$ [is a
  family of subsets of $Y$] generates $\topology_{Y}$ the topology of
$Y$, then $f$ is continuous iff for any $V\in\mathcal{E}$ we have its preimage
$f^{-1}V\in\topology_{X}$ is open.

If you believe these lemmas, then Urysohn's lemma follows immediately.
\end{proof}

\begin{problem}
Given a topological space, can its topology be described by a metric?
\end{problem}

\begin{notation}
We write $(M)$ for the metrizable topological spaces.
\end{notation}

\begin{node}
We know that metrizable spaces are first-countable. And we know
metrizable spaces are normal. But there are first-countable normal
spaces which are not metrizable.
\end{node}

\begin{node}
We can prove that regular Lindel\"{o}f spaces are metrizable, which is
stronger than the result we want to prove today. We want to prove that
second-countable normal spaces are metrizable. This is sufficient, but
there are stronger conditions. There are necessary and sufficient
conditions for metrizable spaces, but they are inelegant.

Note: metrizable spaces are not always second-countable normal spaces.
\end{node}

\begin{node}[Strategy] The proof strategy consists of two steps:
\begin{enumerate}
\item Define \define{Hilbert cube} $H=[0,1]^{\NN}$ the product space
  of countably many intervals; we claim $H$ is metrizable.
\item If $X$ is a second-countable normal space, then there exists a
  map $f\colon X\to H$ which is an emedding (i.e., an injective
  continuous map---equivalently, $f\colon X\to f(X)$ is a
  homeomorphism when we give $f(X)$ the relative topology). Then
  $f(X)$ is metrizable, hence $X$ must be metrizable.
\end{enumerate}
We will need the Baer sets $G_{\delta}$.
\end{node}

\begin{definition}
Let $(X,\topology)$ be a topological space. Let $E\subset X$ be a
subset. We say that $E$ is \define{$G_{\delta}$} if we can represent
$E$ as the countable intersection of open sets.
\end{definition}

\begin{example}
If $X$ is a metric space, then all closed sets are $G_{\delta}$ (just
take all $\varepsilon=1/n$ neighborhoods of the closed set).
\end{example}

\begin{example}
Let $(X,\topology)$ be a topological space. Let $f\colon X\to[0,1]$
be a continuous function. Then $f^{-1}(\{0\})$ the preimage of $0$ is
closed $G_{\delta}$.

This is because $\RR$ is a metric space, therefore $\{0\}$ is a closed
$G_{\delta}$ subset of $\RR$. Then we can lift it using the
preimage. This is because the preimage preserves open sets and
intersections. Therefore the $G_{\delta}$ condition is preserved.
\end{example}

\begin{lemma}
Second-countable normal spaces have closed subsets be $G_{\delta}$.
\end{lemma}
